\chapter{方法与系统设计}

\section{总体流程}

项目流程分为数据层、训练层、评价层和应用层。数据层读取原始目录与标签CSV，生成患者级manifest，并在训练前执行完整性门禁。训练层根据配置创建模型、数据增强、损失函数、优化器和scheduler，以validation loss保存best checkpoint。评价层从checkpoint内嵌设置恢复模型和预处理，输出逐图概率、JSON指标和混淆矩阵。应用层读取冻结权重，为单张上传图像提供与评价一致的推理。

这一设计强调“结果由预测文件生成”。表格中的指标不是手工录入的孤立数字，而能追溯到逐样本路径、真实标签、预测标签和两类概率。数据manifest、配置与checkpoint共同决定一个run的含义。

\section{图像预处理}

所有模型把输入转换为RGB三通道并调整为$224\times224$。Kermany原图纵横比不完全一致，直接resize会产生形变；RSNA处理图多为方形。当前严格基线保持原项目预处理，以便检验数据协议影响。评价阶段不使用随机增强，只执行resize、tensor转换和ImageNet归一化：
\begin{equation}
x'_{c}=\frac{x_c-\mu_c}{\sigma_c},
\end{equation}
其中$\mu=(0.485,0.456,0.406)$，$\sigma=(0.229,0.224,0.225)$。

训练基线采用随机水平翻转和$\pm7^\circ$旋转。稳健方案增加亮度和对比度轻微扰动。胸片左右结构具有一定医学含义，水平翻转可能改变侧别标记；由于任务只做全局二分类、没有使用侧别标签，本项目把其作为基线增强保留，同时在限制中说明。大角度旋转、垂直翻转和强颜色偏移不符合胸片采集规律，未使用。

\section{DenseNet121}

DenseNet的第$l$层接收前面所有层特征拼接：
\begin{equation}
x_l=H_l([x_0,x_1,\ldots,x_{l-1}]).
\end{equation}
密集连接促进低层边缘与高层语义特征复用。项目使用ImageNet预训练DenseNet121，将分类器替换为两输出全连接层。严格设置为AdamW、学习率$10^{-4}$、权重衰减$10^{-4}$、5个epoch和cosine scheduler。DenseNet承担主基线及后续稳健/混合实验，因为其内部表现稳定、训练成本低于ViT。

\section{ConvNeXt-Tiny}

ConvNeXt-Tiny使用现代化卷积stage设计和较大卷积核，在保持卷积归纳偏置的同时吸收Transformer时代的训练经验。项目替换末端分类层为两类输出，使用ImageNet预训练、学习率$5\times10^{-5}$、5个epoch。它的作用是检查“现代卷积结构”是否比DenseNet更能抵抗来源变化，而不是追求模型数量。

\section{ViT-B/16}

ViT-B/16把$224\times224$输入划分为$16\times16$ patch，共196个patch。每个patch线性嵌入后加入位置编码和分类token，通过多层自注意力编码。项目使用ImageNet预训练权重，替换head为两类输出。

审查发现，旧公开配置误写学习率$3\times10^{-4}$和batch 32，而发布checkpoint内嵌真实设置为学习率$5\times10^{-5}$、batch 16、混合精度开启。错误配方在新划分前两epoch的validation准确率仅约0.74和0.78。修正后同seed前两epoch达到约0.89和0.97。正式结果仅使用修正配置；错误run保留日志但不进入汇总。

\section{损失函数和优化}

基线使用按训练集类别频数计算的加权交叉熵。若类别$c$有$n_c$个样本、总数为$N$、类别数为$C$，权重设为
\begin{equation}
w_c=\frac{N}{Cn_c}.
\end{equation}
稳健方案加入标签平滑$\epsilon=0.1$，目标分布从one-hot变为
\begin{equation}
\tilde y_c=(1-\epsilon)y_c+\epsilon/C.
\end{equation}
它可以抑制极端logit，但也可能改变固定阈值概率，因此必须联合观察AUC、Brier和ECE。

AdamW把权重衰减与梯度更新分离。cosine scheduler在预设epoch内逐步降低学习率。checkpoint按validation loss最小保存，不按test结果选择。当前训练没有在每个epoch提前停止，以保证各seed预算一致；best checkpoint仍可来自不同epoch。

\section{随机性控制}

训练同时设置Python、NumPy、CPU和CUDA随机种子，关闭cuDNN benchmark并请求确定性算法。RTX 5080上的Flash/Memory-Efficient Attention会报告部分反向算子无法完全确定，程序以warn-only记录。为避免把“设置了seed”误写成“位级完全可复现”，本文把三seed统计作为主要稳定性证据，并保留环境和警告日志。

\section{混合训练与域均衡}

患者级Kermany train/validation与RSNA对应集合合并，文件名前加\texttt{kermany\_}或\texttt{rsna\_}保存来源。简单混合训练按全部图像普通shuffle，因此样本更多的Kermany占更大权重。域均衡采样为来源$d$中的每个样本赋予$1/n_d$权重，再有放回抽样，使两个来源期望采样概率近似相同。

域均衡只平衡来源，不同时保证来源内类别平衡；类别不均衡仍由损失权重处理。这种两层机制可能增加梯度方差，因此必须用多seed检查。混合训练使用同一DenseNet结构和5个epoch，内部与RSNA测试分别评价。

\section{概率校准与阈值}

模型logit除以温度$T>0$后再softmax：
\begin{equation}
p_T(y=k\mid x)=\frac{\exp(z_k/T)}{\sum_j\exp(z_j/T)}.
\end{equation}
$T$只能在validation预测上按NLL拟合。$T>1$通常降低极端置信，$T<1$增强置信。内部test使用Kermany validation拟合温度；RSNA test若做目标域校准，只能使用RSNA validation。两种情形回答的问题不同：前者是同来源内部校准，后者是在目标域已有标注校准数据时的适配。

默认阈值0.5便于比较。若选择冻结操作点，则在validation上最大化预设指标或满足召回下限后固定，再应用test。test阈值扫描仅画权衡曲线，不报告为部署最优阈值。

\section{错误分析与Grad-CAM}

错误分析把样本划分为FP、FN、高置信错误和概率0.4--0.6边界样本。FP反映阴性误报，FN反映肺炎漏检。高置信错误门槛预先设定，例如错误类别概率不低于0.9。错误网格按固定规则选取概率最极端与最接近阈值的样本，避免只挑视觉上易解释的图。

Grad-CAM对DenseNet最后特征层计算目标类别梯度并加权汇总。本文展示成功与失败案例，解释模型使用的区域，但不把热图当作病灶分割。RSNA框只可用于粗粒度一致性探索。

\section{网页演示}

网页由浏览器前端和Flask推理接口组成。后端启动时加载checkpoint和类别顺序，上传后检查文件、读取PIL图像、执行checkpoint内嵌预处理并返回JSON概率。页面展示NORMAL和PNEUMONIA概率、阈值、预测倾向和明确警示。网页不在线训练，不保存患者资料，也不允许根据上传图像修改阈值。

\section{参数量、计算量与比较公平性}

DenseNet121、ConvNeXt-Tiny与ViT-B/16的参数规模和计算图不同。ViT-B/16参数更多、batch受16GB显存限制为16；两个CNN可用batch 32。若强制完全相同batch，可能让ViT无法运行；若各自使用合理batch，又不能把差异解释为纯架构因果效应。

本文采用“实用配方公平性”：每个模型使用已验证可训练的预训练设置，epoch接近，checkpoint按同一validation loss规则选择，每个设置运行相同seed数。它回答的是课程项目中三条可实施路线的整体表现，而不是控制所有训练变量后的架构消融。

为了减少选择机会差异，未对某个模型单独进行大量test驱动搜索。ViT学习率修正来自发布checkpoint内嵌真实记录，而不是根据新test挑选。DenseNet作为后续方法主干，是在内部稳定性、训练成本和旧胸片使用经验基础上选择，不代表其test分数最高。

\section{checkpoint选择与过拟合}

每个epoch结束后计算validation加权交叉熵，若低于历史最优便保存best checkpoint，同时始终保存last checkpoint。报告评价best而非last。validation loss同时受类别权重和概率尺度影响，可能选择accuracy相近但置信不同的模型。

训练只运行4--5个epoch，没有长时间训练后再从大量epoch挑选，因此选择机会有限。曲线显示多数CNN在2--5 epoch已经接近饱和，后期train loss继续下降而validation loss波动，提示进一步训练可能过拟合。ViT正确学习率下也在4 epoch内达到约0.97 validation accuracy。

严格的后续工作可以预设patience并最多训练20 epoch，但不同run实际epoch会变化。当前固定短预算更适合课程比较，也降低算力消耗。

\section{域策略的目标函数解释}

简单混合最小化所有训练图像平均损失：
\begin{equation}
\mathcal{L}_{mix}=\frac{1}{N_K+N_R}\left(\sum_{i\in K}\ell_i+\sum_{j\in R}\ell_j\right).
\end{equation}
由于$N_K>N_R$，Kermany贡献更大。域均衡近似最小化两个域平均损失的均值：
\begin{equation}
\mathcal{L}_{db}=\frac{1}{2}\left(\frac{1}{N_K}\sum_{i\in K}\ell_i+\frac{1}{N_R}\sum_{j\in R}\ell_j\right).
\end{equation}
实现通过WeightedRandomSampler近似，而不是在每个batch显式包含等量两域。某些batch仍可能不平衡，有放回抽样也会重复样本。这些实现细节解释了域均衡的随机波动。

GroupDRO会进一步优化最大组损失，CORAL会对齐特征协方差。它们需要明确组、对齐层与权重，并增加超参数。简单混合达到预设门槛后，本项目按停止规则没有在同一RSNA test继续搜索这些方法。
